\documentclass[11pt,a4paper]{article}
\usepackage[margin=25mm]{geometry}
\usepackage{amsmath,amssymb,amsthm,mathtools,booktabs,microtype}
\usepackage[colorlinks=true,urlcolor=blue,citecolor=blue,linkcolor=blue]{hyperref}
\newtheorem{theorem}{Theorem}
\newtheorem{proposition}[theorem]{Proposition}
\newtheorem{corollary}[theorem]{Corollary}
\newtheorem{lemma}[theorem]{Lemma}
\newcommand{\E}{\mathbb E}
\newcommand{\R}{\mathbb R}
\newcommand{\diag}{\operatorname{diag}}
\newcommand{\Tr}{\operatorname{Tr}}
\title{A sharp two-rate reset witness against static qutrit disorder}
\author{Lluis Eriksson\\\small Research manuscript prepared with AI assistance}
\date{12 October 2026 --- research candidate, not deposited}
\begin{document}
\maketitle
\begin{abstract}
Can stationary states at a few reset rates distinguish fluctuating
dephasing from an arbitrary distribution of static Hamiltonians?
For a qubit, every finite list of stationary states of a real dephasing
semigroup has a finite static alias. We exhibit a sharp obstruction for
a qutrit. If $c_{jk}(\nu)$ is the real, normalized stationary coherence
under Poisson resets and $A(\nu)=[c_{12}(\nu)+c_{23}(\nu)]/2$, every
static diagonal Hamiltonian ensemble, without support or moment
restrictions, satisfies
\[
 A(a)+9A(4a)-5c_{13}(4a)\leq5\qquad(a>0).
\]
The completely positive semigroup with generator
$-\gamma[D,[D,\cdot]]$, $D=\diag(0,1,2)$, gives $16/3$ at
$a=2\gamma$. We prove the inequality by an explicit polynomial
certificate, determine all its finite equality points, and show that the
exact infimum of the uniform error in the three selected statistics is
$1/45$. The infimum is approached by seven-point static ensembles but
is attained by no probability measure on $\R^2$. A single reset rate
cannot exclude static disorder for this semigroup. We also give a
deterministic error rule, a finite-shot implementation, and a spectral
coupling that extends the exclusion to arbitrary static qutrit Hamiltonian
ensembles with noisy population calibration. If all three population
leakages at rate $a$ are at most $\delta$, the bound becomes
$W_a\leq5+30\sqrt{2\delta}+30\delta$. This is a witness against static
disorder, not against classical time-dependent noise or a proof of a
quantum environment.
\end{abstract}

\section{Question and relation to earlier work}
Static Hamiltonian ensembles are an established description of
inhomogeneous broadening. In particular, Cauchy-distributed qubit
frequencies reproduce exponential dephasing exactly
\cite{kropf,chen2018}. Higher-dimensional dephasing imposes additional
consistency: three transition frequencies of a qutrit must be
differences of the same three energies. Chen et al.\ explicitly formulate
this joint-distribution constraint using the root system of
$\mathfrak{su}(3)$ \cite{chen2019}. The distinction between Hamiltonian
ensemble representability and other notions of classicality or
Markovianity is also established \cite{chen2021}.
Pulse-based discrimination of static broadening is much older: the Hahn
echo refocuses static frequency disorder already for a qubit \cite{hahn}.
Our restrictions are free evolution between preparations and stationary
or exponentially sampled observations, without echo or decoupling pulses.
We claim no advantage over those controlled protocols.
Classical stochastic dilations of dephasing semigroups are studied
systematically in \cite{kummerer,fagnola}. Finite-data certification of different
environmental properties, including tests based on failure of
mixed-unitarity, is already available \cite{lin}; our selected statistics
do not certify that stronger property.

The question here is narrower and operational: can \emph{two stationary
reset settings}, without reconstructing a full time-dependent channel,
exclude every static diagonal ensemble? We give an explicit yes answer,
an exact bounded-observable inequality, and a sharp distance to the
allowed data. The conceptual incompatibility of qutrit marginals is
prior work. The candidate contribution is the finite-resolvent witness
and its sharp error geometry. The separating-hyperplane method is
classical generalized-moment geometry \cite{kemperman}, not a new
general method of moment inference. The literature comparison does not
establish that nobody has previously derived the same inequality.

Stationary reset states are Laplace transforms of free evolution; this
renewal description is standard \cite{rose}. Their diagnostic value
depends on which controls and preparations are held fixed. The same
probability law of static Hamiltonians must be used at all settings.
Allowing a different law at each reset rate destroys the proposed test.
The author's earlier discussions of maintenance and restoration
\cite{eriksson-maintenance} and finite-data moment falsification
\cite{eriksson-moments} motivated this question. These unreviewed
manuscripts are contextual only; no theorem from them is used here,
and no thermodynamic work claim is made.

\section{Protocol and the model being excluded}
Set $\hbar=1$. In each realization draw $(X,Y)$ from an arbitrary Borel
probability measure $\mu$ on $\R^2$ and hold
\begin{equation}\label{eq:H}
 H_{X,Y}=\diag(0,X,X+Y)
\end{equation}
fixed. No independence, symmetry, density, finite-support, bandwidth or
moment assumption is imposed on $\mu$.
Reset the system, independently of its sampled energies, to
$\rho_+=|+_3\rangle\langle+_3|$, where
$|+_3\rangle=(|1\rangle+|2\rangle+|3\rangle)/\sqrt3$, at Poisson
rate $\nu>0$. The static register is not reset. After taking the
stationary limit, average across independently prepared realizations.
This is ensemble tomography, not an ergodic time average over a single
conserved register.

For a signed energy gap $Z$, the normalized stationary coherence is
\begin{equation}\label{eq:renewal}
 z_Z(\nu)=\int_0^\infty\nu e^{-\nu t}\E e^{iZt}\,dt
 =\E\frac{\nu}{\nu-iZ},\qquad
 c_Z(\nu):=\Re z_Z(\nu)=\E g_\nu(Z),
 \quad g_\nu(z)=\frac{\nu^2}{\nu^2+z^2}.
\end{equation}
Boundedness justifies the interchange of expectations and integrals.
Operationally this is also a Laplace-weighted free-induction-decay
(FID/Ramsey) measurement: prepare independently, draw a waiting time
from $\operatorname{Exp}(\nu)$, and measure once. Physical repeated
resetting is therefore not necessary to obtain the same averaged data.
The preparation and Hamiltonian law must still be the same at every
setting; compare the FID formulation in \cite{chen2018}.
Thus $c_{jk}(\nu)=3\Re(\rho_{\mu,\nu})_{jk}$ and
\begin{equation}\label{eq:features}
 c_{12}(\nu)=\E g_\nu(X),\quad
 c_{23}(\nu)=\E g_\nu(Y),\quad
 c_{13}(\nu)=\E g_\nu(X+Y),\quad
 A(\nu)=\frac{c_{12}(\nu)+c_{23}(\nu)}2.
\end{equation}
All populations equal $1/3$. Imaginary coherences are allowed in the
excluded model; the witness uses their real parts only.

The reference process is the legitimate Lindblad semigroup
\begin{equation}\label{eq:ref}
 \mathcal L_\gamma(\rho)=-\gamma[D,[D,\rho]],\quad
 D=\diag(0,1,2),\qquad\gamma>0.
\end{equation}
It has jump operator $\sqrt{2\gamma}D$. Its adjacent coherences decay
as $e^{-\gamma t}$ and the outer coherence as $e^{-4\gamma t}$.
Under the same reset protocol,
\begin{equation}\label{eq:refdata}
 A_M(\nu)=\frac\nu{\nu+\gamma},\qquad
 c_{13,M}(\nu)=\frac\nu{\nu+4\gamma}.
\end{equation}
This semigroup also has a classical \emph{time-dependent} Gaussian
noise realization: the random phase $\sqrt{2\gamma}W_tD$, with
standard Brownian motion $W_t$, produces precisely these factors.
It follows that rejecting static ensembles does not certify quantum
noise, nonclassical dynamics in every usage of that term, or
non-Markovianity.

\section{The sharp universal inequality}
\begin{theorem}[Two-rate static-disorder bound]\label{thm:witness}
For every probability measure $\mu$ in \eqref{eq:H} and every $a>0$,
\begin{equation}\label{eq:witness}
 W_a:=A(a)+9A(4a)-5c_{13}(4a)\leq5.
\end{equation}
The constant $5$ is attained. Equality holds exactly when $\mu$ is
supported on
\begin{equation}\label{eq:equality}
 (0,0),\qquad(a\sqrt2,a\sqrt2),\qquad(-a\sqrt2,-a\sqrt2).
\end{equation}
\end{theorem}

\begin{proof}
Rescale $x=X/a$, $y=Y/a$ and write
$A_r(x,y)=[g_r(x)+g_r(y)]/2$. It is enough to show
\begin{equation}\label{eq:pointwise}
 h(x,y):=5-A_1(x,y)-9A_4(x,y)+5g_4(x+y)\geq0.
\end{equation}
Let
\[
 \Delta=(1+x^2)(16+x^2)(1+y^2)(16+y^2)(16+(x+y)^2)>0.
\]
Clearing denominators gives the exact identity
\begin{equation}\label{eq:certificate}
 h(x,y)=\frac{5P(x,y)}{2\Delta}.
\end{equation}
For $p=xy$ and $d=(x-y)^2$, the polynomial has the form
\begin{align}\label{eq:positivecase}
 P={}&(5p^2+48)d^2+2B(p)d
       +8p(p-2)^2(p+1)(p+16),\\
 B(p)={}&p^4+15p^3+108\left(p-\frac{46}{27}\right)^2
                +\frac{1904}{27}.\nonumber
\end{align}
If $p\geq0$, all terms are nonnegative and $B(p)>0$.
Equality requires $d=0$ and $p\in\{0,2\}$, giving the three points
in \eqref{eq:equality} after rescaling.

If $p<0$, set $q=-p>0$ and $s=(x+y)^2\geq0$. The same polynomial is
\begin{align}\label{eq:negativecase}
 P={}&(5q^2+48)s^2
 +(2q^4+10q^3+216q^2+1120q+768)s\\
 &+64q^4+1248q^3+3744q^2+2560q>0.\nonumber
\end{align}
This proves the pointwise inequality and all its zeros. Integrate
\eqref{eq:pointwise} against $\mu$. Since $h$ is bounded,
nonnegative, and vanishes only at the stated points, its integral is
zero if and only if the measure is supported there. No moment
condition has entered the argument.
\end{proof}
The equality set above consists of finite points. There is also
asymptotic saturation along the coordinate axes:
\[
 h(x,0)=h(0,x)=\frac{15x^2}{2(1+x^2)(16+x^2)}\longrightarrow0.
\]
These escaping directions explain the limiting construction in
Theorem~\ref{thm:distance}; they are not additional points in $\R^2$.

\begin{corollary}[Violation by a completely positive semigroup]
For \eqref{eq:ref}, put $t=\gamma/a$. Then
\begin{equation}\label{eq:violation}
 W_{a,M}-5=\frac{t(7-5t)}{(1+t)(4+t)}.
\end{equation}
It is positive for $0<t<7/5$, and its maximum over $a>0$ is $1/3$,
attained at $a=2\gamma$. In particular,
\[
 A_M(2\gamma)=\frac23,\qquad A_M(8\gamma)=\frac89,
 \qquad c_{13,M}(8\gamma)=\frac23,
 \qquad W_{2\gamma,M}=\frac{16}{3}.
\]
\end{corollary}
\begin{proof}
Substitute \eqref{eq:refdata}. The derivative of the right side of
\eqref{eq:violation} is
$-4(2t-1)(4t+7)/[(t+1)^2(t+4)^2]$.
\end{proof}
The optimizing statement concerns this particular witness family with
rate ratio four, not all possible witnesses or reset protocols.

\section{An exact error floor and escape to infinity}
Use the three selected real statistics, in scaled coordinates,
\begin{equation}\label{eq:selected}
 T(\mu)=\left(\E A_1(X,Y),\E A_4(X,Y),\E g_4(X+Y)\right),
 \qquad b=\left(\frac23,\frac89,\frac23\right).
\end{equation}
Here $X,Y$ are dimensionless; physical gaps are $aX,aY$ with
$a=2\gamma$. The vector in \eqref{eq:selected} deliberately does
\emph{not} include $c_{13}(a)$ or all entries of the two stationary
states.

\begin{theorem}[Sharp distance for the selected statistics]\label{thm:distance}
Over all Borel probability measures on $\R^2$,
\begin{equation}\label{eq:distance}
 \inf_\mu\|T(\mu)-b\|_\infty=\frac1{45}.
\end{equation}
No measure attains the infimum. It is approached by symmetric
seven-point ensembles with unbounded support radius.
\end{theorem}
\begin{proof}
The linear functional $\ell(u)=u_1+9u_2-5u_3$ has dual norm
$\|\ell\|_1=15$, and $\ell(b)=16/3$. Theorem~\ref{thm:witness}
therefore gives $15\|T(\mu)-b\|_\infty\geq1/3$.

For the matching upper bound, assign mass $19/45$ to $(0,0)$, mass
$1/5$ to each of $(\sqrt2,\sqrt2)$ and $(-\sqrt2,-\sqrt2)$, and
mass $2/45$ to each of $(0,R),(R,0),(0,-R),(-R,0)$, where $R>0$.
These seven masses sum to one. Their statistics are
\begin{equation}\label{eq:approx}
 T(\mu_R)=
 \left(\frac{29}{45}+\frac4{45}g_1(R),\quad
       \frac{13}{15}+\frac4{45}g_4(R),\quad
       \frac{31}{45}+\frac8{45}g_4(R)\right).
\end{equation}
Consequently
\begin{equation}\label{eq:convergence}
 \|T(\mu_R)-b\|_\infty
 =\frac1{45}+\frac{128}{45(16+R^2)}\longrightarrow\frac1{45}.
\end{equation}
For the equality obstruction, attaining $1/45$ forces equality in
both the linear bound and the static witness. All three nonzero
coefficients then force
$T=(29/45,13/15,31/45)$. But equality in the static witness requires
support on \eqref{eq:equality}; every such measure has its second
coordinate at least $8/9$, whereas $13/15<8/9$. This is impossible.
\end{proof}

Seven points are an explicit symmetric construction, not an optimal
atom count. For the three selected real statistics alone, every
attainable vector has a representation with at most four atoms:
apply the measurable-function version of Tchakaloff's theorem
\cite[Corollary 2]{bayer} to $1,A_1,A_4,g_4(X+Y)$.
Such a reduction need not preserve the other or imaginary coherences.

Thus the set of attainable selected data is not closed. This point
matters in numerical optimization: an apparent optimum can represent
probability mass moving to arbitrarily large energy gaps. On any fixed
compact gap domain the probability measures form a compact set and
$T$ is continuous, so the best error is attained and is strictly larger
than $1/45$. The theorem gives the infimum for the three selected
statistics, not a sharp trace-distance error for the full states.

\section{Why the second rate and the third level matter}
\begin{proposition}[One rate is insufficient for the reference example]
For every single rate $\nu>0$, a two-point static qutrit ensemble
reproduces the entire stationary state of \eqref{eq:ref} exactly.
Indeed, take $H=\pm\sqrt{\gamma\nu}D$ with equal probability.
\end{proposition}
\begin{proof}
For a gap $k\sqrt{\gamma\nu}$ with $k=1,2$, the averaged complex
coherence is real and equals
$g_\nu(k\sqrt{\gamma\nu})=\nu/(\nu+k^2\gamma)$.
These are exactly the adjacent and outer coherences; populations
agree as well.
\end{proof}

\begin{proposition}[Finite qubit aliases]\label{prop:qubit}
For a qubit with real Markovian dephasing rate $\gamma>0$ and any
$N$ distinct positive reset rates, a symmetric finite static frequency
distribution reproduces all stationary states exactly. One can use
$N+1$ signed frequencies, and fewer are impossible.
\end{proposition}
\begin{proof}
Let $\nu_1,\ldots,\nu_N$ be the rates and write
$Q(s)=\prod_j(s+\nu_j)=E(s)+O(s)$, separating even and odd powers.
Set
\begin{equation}\label{eq:qubitF}
 F(s)=
 \begin{cases}
 O(s)/(sO(s)+\gamma E(s)),&N\text{ odd},\\
 E(s)/(sE(s)+\gamma O(s)),&N\text{ even}.
 \end{cases}
\end{equation}
Since $Q(-\nu_j)=0$, $E(\nu_j)=O(\nu_j)$, giving
$F(\nu_j)=1/(\nu_j+\gamma)$.

The denominator is the degree-$N+1$ parity part of
$(s+\gamma)Q(s)$. For $\omega\geq0$ the phase
\[
 \theta(\omega)=\arctan(\omega/\gamma)
                  +\sum_j\arctan(\omega/\nu_j)
\]
is strictly increasing from $0$ to $(N+1)\pi/2$. Its crossings of
$(k+1/2)\pi$ when $N$ is odd, or $k\pi$ when $N$ is even,
give all the nonnegative imaginary-axis poles. The poles are simple,
come in pairs, and include zero exactly when $N$ is even. Their total
number is $N+1$. Each nonzero signed pole has residue
\begin{equation}\label{eq:mass}
 p_\omega=\frac{\gamma}{(\gamma^2+\omega^2)
 [\gamma/(\gamma^2+\omega^2)+\sum_j\nu_j/(\nu_j^2+\omega^2)]}>0.
\end{equation}
When present, the zero residue is
$p_0=(1+\gamma\sum_j\nu_j^{-1})^{-1}>0$.
Also $sF(s)\to1$, so the residues sum to one. Partial fractions
therefore express $F$ as the Cauchy transform of a symmetric
probability measure on these frequencies.

For the lower bound let $v_j(\omega)=(\nu_j+i\omega)^{-1}$ and
$f_j=(\nu_j+\gamma)^{-1}$. Any representing measure has the same
Gram matrix of $1,v_1,\ldots,v_N$, whose lower-right block is
$K_{ij}=(f_i+f_j)/(\nu_i+\nu_j)$. Its Schur complement is
\[
 K-ff^T=2\gamma\,\diag(f)C\diag(f),\qquad
 C_{ij}=\frac1{\nu_i+\nu_j}
       =\int_0^\infty e^{-\nu_i t}e^{-\nu_j t}\,dt>0.
\]
The distinct exponentials are linearly independent, so the Gram
matrix has rank $N+1$. A measure with fewer than $N+1$ atoms
cannot have that rank.
\end{proof}

The construction in Proposition~\ref{prop:qubit} is a specialization
of classical positive rational quadrature. In particular Lubinsky's
reciprocal-polynomial reproducing kernel \cite[equations (1.11)--(1.12)]{lubinsky},
for weight proportional to
$[(\gamma^2+\omega^2)\prod_j(\nu_j^2+\omega^2)]^{-1}$, produces
the same phase nodes and, after multiplying by
$\prod_j(\nu_j^2+\omega^2)$, the masses \eqref{eq:mass}.
We do not claim this quadrature formula as a new discovery; the proof
is included to make the qubit/qutrit contrast self-contained.
Positive Cauchy transforms and their Gram-matrix constraints also
belong to classical Nevanlinna--Pick interpolation; see
\cite[Theorem 1.2 and the Herglotz representation]{agler}.

The minimality statement is consequently precise: for real
dephasing-semigroup reset data with a fixed dephasing basis, no qubit
provides this finite-rate rejection, while the displayed qutrit does;
and this particular qutrit cannot be rejected at one rate. It is not
a minimality claim over arbitrary controls, echoes, preparations or
general non-dephasing processes.

\section{Compatibility of the full qutrit semigroup}
The finite witness can be related to a simpler full-trajectory fact.
The latter is a Cauchy-marginal compatibility problem, within the
classical joint-mixability framework \cite{wang}; it is not presented
as the main novelty of this paper.
In particular, the sufficiency follows from Wang and Wang's
Corollary 3.6 for unimodal symmetric location-scale families. We
include an elementary construction and a fractional-moment necessity
argument, without applying a finite-mean assumption to Cauchy laws.
Mixability without finite means, including Cauchy laws, has also been
studied directly in \cite{puccetti}; absence of a mean is not in itself
a novel difficulty introduced by this example.

\begin{proposition}[Three Cauchy marginals]
For nonnegative real dephasing rates
$\Gamma_{12},\Gamma_{23},\Gamma_{13}$, a static diagonal ensemble
reproduces all factors $e^{-\Gamma_{jk}t}$ for every $t\geq0$ if
and only if the three rates satisfy all triangle inequalities.
\end{proposition}
\begin{proof}
The characteristic functions, extended to negative times by conjugacy,
force $X,Y,X+Y$ to have centered Cauchy laws of scales
$\Gamma_{12},\Gamma_{23},\Gamma_{13}$, with scale zero interpreted
as a point mass. For $0<p<1$,
\[
 |X+Y|^p\leq|X|^p+|Y|^p,
 \qquad \E|\Gamma Z|^p=\Gamma^p\E|Z|^p<\infty
\]
for a standard Cauchy $Z$. Taking expectations and then $p\uparrow1$
gives $\Gamma_{13}\leq\Gamma_{12}+\Gamma_{23}$. The other two
inequalities follow by writing $X=(X+Y)-Y$ and $Y=(X+Y)-X$.

Conversely, let $E_1,E_2,E_3$ be independent centered Cauchy
variables with scales
$b_j=(\Gamma_{jk}+\Gamma_{j\ell}-\Gamma_{k\ell})/2\geq0$.
Then $E_k-E_j$ is centered Cauchy of scale
$b_j+b_k=\Gamma_{jk}$. The ensemble $\diag(E_1,E_2,E_3)$,
after removing the irrelevant common energy, is of the form
\eqref{eq:H} and gives all the desired factors.
\end{proof}

The reference rates $(\gamma,\gamma,4\gamma)$ violate this
condition despite generating a completely positive semigroup.
For comparison, complete positivity of real qutrit dephasing requires
triangle inequalities for the square roots of the rates
\cite{schirmer}, a different condition.
Theorem~\ref{thm:witness} replaces full-trajectory information by
two reset settings and bounded observables for this example.

\section{Error control and measurement}
Suppose estimates of the selected statistics obey simultaneous
deterministic error bounds
\[
 |\widehat A(a)-A(a)|\leq\varepsilon_1,\quad
 |\widehat A(4a)-A(4a)|\leq\varepsilon_2,\quad
 |\widehat c_{13}(4a)-c_{13}(4a)|\leq\varepsilon_3.
\]
Every static diagonal model must then satisfy
\begin{equation}\label{eq:errors}
 \widehat W_a\leq5+\varepsilon_1+9\varepsilon_2+5\varepsilon_3.
\end{equation}
This is a model-exclusion rule for a given data set. In particular,
the ideal reference vector cannot be matched to within uniform
error $\varepsilon<1/45$. If the \emph{estimates themselves} may
deviate from the ideal reference by $\varepsilon$, and the same
size error bar is used in \eqref{eq:errors}, a conservative
worst-case guarantee of rejection requires $\varepsilon<1/90$.
These two different error statements must not be conflated.

For a direct finite-shot implementation, measure
\[
 M_{jk}=\tfrac32(|j\rangle\langle k|+|k\rangle\langle j|).
\]
Its expectation equals $c_{jk}$ and its outcomes lie in
$\{-3/2,0,3/2\}$. Use independent stationary preparations for the
five settings $(12,a),(23,a),(12,4a),(23,4a),(13,4a)$, with $m$
shots per setting. A union bound and Hoeffding's inequality give
simultaneous errors at most $e$ with probability at least $1-\alpha_c$
whenever
\begin{equation}\label{eq:shots}
 m\geq\frac{9}{2e^2}\log\frac{10}{\alpha_c}.
\end{equation}
The averages defining $A$ then also have error at most $e$.
This conservative bound is not an optimal experimental allocation.
Independence here means independently prepared static registers, not
repeated correlated measurements on a single conserved register.

\section{Robust population calibration for general static ensembles}
The sharp witness and distance concern diagonal Hamiltonians. We
next give a quantitative reduction from an arbitrary static ensemble
to one diagonal ensemble that works simultaneously at every reset rate.
For any law $\mu$ on Hermitian matrices, define its reset channel by
\[
 \mathcal R_\nu(\rho)=\E_\mu\int_0^\infty\nu e^{-\nu t}
                       e^{-iHt}\rho e^{iHt}\,dt.
\]

\begin{theorem}[A common diagonal surrogate]\label{thm:surrogate}
Let $\mu$ be any probability law on finite Hermitian $d\times d$
matrices, and let $\mathcal R_\nu$ be its stationary reset channel.
Define
\[
 L_j(\nu)=1-\langle j|\mathcal R_\nu(P_j)|j\rangle,
 \qquad P_j=|j\rangle\langle j|.
\]
There is one probability law $\widetilde\mu$ on diagonal Hamiltonians,
independent of $\nu$, such that the normalized coherences from the
uniform reset state satisfy, for all $j\ne k$ and all $\nu>0$,
\begin{equation}\label{eq:surrogate}
 |z_{jk,\mu}(\nu)-z_{jk,\widetilde\mu}(\nu)|
 \leq \sqrt{d-1}\bigl(\sqrt{L_j(\nu)}+\sqrt{L_k(\nu)}\bigr)
 +(d-1)\sqrt{L_j(\nu)L_k(\nu)}.
\end{equation}
Here $z_{jk,\mu}(\nu)=d\langle j|\mathcal R_\nu(\rho_+)|k\rangle$
and $\rho_+$ is the uniform pure state. No norm or moment bound on
the Hamiltonian distribution is needed.
\end{theorem}
\begin{proof}
Conditional on $H$, let $\eta_j^H(B)=\langle j|P_H(B)|j\rangle$
be its local spectral probability measure. Draw energies
$E_1,\ldots,E_d$ independently conditional on $H$, with respective
laws $\eta_j^H$, and set $\widetilde H=\diag(E_1,\ldots,E_d)$.
The resulting diagonal Hamiltonian law is independent of time and
reset rate. This construction uses spectral measures and does not
require an arbitrary eigenbasis choice at degenerate eigenvalues.
For $U=e^{-iHt}$, its conditional free coherence is
$U_{jj}\overline{U_{kk}}$.

For the original uniform pure input, the normalized free coherence is
$(U_{jj}+s_j)\overline{(U_{kk}+s_k)}$, where
$s_j=\sum_{\ell\ne j}U_{j\ell}$. Put
$l_j=1-|U_{jj}|^2\in[0,1]$. Unitarity and Cauchy--Schwarz give
$|s_j|\leq\sqrt{d-1}\sqrt{l_j}$. The same $l_j$ equals the
total probability of leaving $|j\rangle$, since the squared norm
of both row $j$ and column $j$ is one.
Since $|U_{jj}|\leq1$, the absolute difference of the two coherences
is bounded by
\[
 \sqrt{d-1}(\sqrt{l_j}+\sqrt{l_k})+(d-1)\sqrt{l_jl_k}.
\]
Average over $H$ and over the exponential age distribution
$\nu e^{-\nu t}dt$. This is a probability space, and
$\E l_j=L_j(\nu)$. Jensen's and Cauchy--Schwarz inequalities
give \eqref{eq:surrogate}. All integrands are bounded, so no moment
or dominated-differentiation assumption is necessary.
\end{proof}

\begin{corollary}[Noisy population-calibrated qutrit witness]\label{cor:robust}
For any static qutrit Hamiltonian ensemble, if
$L_j(a)\leq\delta$ for all $j=1,2,3$, then its uniform-reset data obey
\begin{equation}\label{eq:robust}
 W_a\leq5+30\sqrt{2\delta}+30\delta.
\end{equation}
Consequently, with simultaneous coherence error bounds as in
\eqref{eq:errors} and simultaneous certified population bounds
$L_j(a)\leq\delta$, observing
\[
 \widehat W_a>5+30\sqrt{2\delta}+30\delta
              +\varepsilon_1+9\varepsilon_2+5\varepsilon_3
\]
excludes every static qutrit Hamiltonian ensemble using the same
Hamiltonian law for all the tested reset rates and preparations.
\end{corollary}
\begin{proof}
For a fixed Hamiltonian, the population leakage has the positive
spectral representation
\begin{equation}\label{eq:monotone}
 L_j(\nu)=\E_H\iint
 \frac{(E-E')^2}{\nu^2+(E-E')^2}\,
 d\eta_j^H(E)\,d\eta_j^H(E').
\end{equation}
Indeed, the return amplitude is $U_{jj}(t)=\int e^{-iEt}d\eta_j^H(E)$;
expanding its squared modulus and taking the exponential age average
gives \eqref{eq:monotone}. Thus $L_j(\nu)$ is nonincreasing in
$\nu$, without any moment restriction. In particular
$L_j(4a)\leq L_j(a)\leq\delta$.
Theorem~\ref{thm:surrogate} bounds each coherence error relative to
its one common diagonal surrogate by $2\sqrt{2\delta}+2\delta$
at both rates. The same bound holds for the adjacent averages.
Apply Theorem~\ref{thm:witness} to the surrogate and sum the errors:
\[
 (1+9+5)(2\sqrt{2\delta}+2\delta)
 =30\sqrt{2\delta}+30\delta.
\]
\end{proof}

The constants in \eqref{eq:robust} are conservative, not claimed
optimal. The ideal reference has $L_j=0$. In the absence of coherence
error, its margin $1/3$ exceeds the correction whenever
\[
 \delta<\left(\frac{2\sqrt{115}-15\sqrt2}{30}\right)^2.
\]
For finite data, population upper bounds must have their own stated
joint confidence. Sampling a fresh register for each shot remains
essential; an independent-shot analysis cannot be applied to one
uncontrolled conserved register. The sharp $1/45$ distance from
Theorem~\ref{thm:distance} remains a statement about the diagonal
model and selected statistics, not a sharp error for this enlarged test.

The practical requirements of this conservative bound are demanding.
With no coherence error, the reference is certified only when
$\delta<\delta_*\simeq6.1052\times10^{-5}$. With five simultaneous
coherence errors at most $e$, the worst-case reference guarantee is
\[
 30e+30\sqrt{2\delta}+30\delta<\frac13.
\]
For example, $e=1/180$ and $\alpha_c=0.025$ require about
$8.74\times10^5$ shots \emph{per coherence setting} under
\eqref{eq:shots}. For each of the three population preparations,
$n$ independent shots with zero escapes give the one-sided upper
bound $L_j\leq\log(3/\alpha_p)/n$, simultaneously except with
probability at most $\alpha_p$: for any larger leakage, the
probability of zero escapes is at most $e^{-nL_j}$, and a union
bound covers the three states. Taking $\alpha_p=0.025$ requires
about $4.79\times10^5$ zero-escape shots per state to certify
$\delta\leq10^{-5}$. The combined coverage is at least
$1-\alpha_c-\alpha_p$. If escapes occur, a binomial confidence
bound must replace the zero-event formula. These estimates account
for sampling only, not state-preparation or measurement bias, and
are not evidence that an efficient experiment has been achieved.

Exact data admit a simpler calibration with only two basis states.

\begin{proposition}[Exact population calibration]\label{prop:diag}
Let $\mu$ be a probability measure on arbitrary finite Hermitian
$3\times3$ matrices, with $\mathcal R_\nu$ as defined above.
If, at any one $\nu>0$,
$\mathcal R_\nu(P_1)=P_1$ and $\mathcal R_\nu(P_2)=P_2$,
where $P_j=|j\rangle\langle j|$, then $H$ is diagonal almost surely.
Consequently, these two exact calibrations together with
$W_a>5$ exclude every static qutrit Hamiltonian ensemble with the
same Hamiltonian law for all reset preparations and rates.
\end{proposition}
\begin{proof}
For each $j=1,2$ and $k\ne j$, the assumption gives
\[
 0=\E_\mu\int_0^\infty\nu e^{-\nu t}
                 |\langle k|e^{-iHt}|j\rangle|^2\,dt.
\]
The integrand is nonnegative. For almost every $H$, all these
finitely many amplitudes vanish for almost every $t$, and hence for
every $t$ by continuity. Differentiating at zero gives
$H_{kj}=0$. Hermiticity and dimension three leave no off-diagonal
entry.
\end{proof}

The reference semigroup passes these population calibrations exactly.
Neither the exact nor the noisy calibration excludes arbitrary
system--bath Hamiltonians with a quantum bath or preparation-dependent
changes in the Hamiltonian distribution.

\section{Verification, scope and open questions}
The accompanying verifier checks the rational identity
\eqref{eq:certificate}, both polynomial representations, the
equality points, the semigroup violation and its optimizing rate,
and the exact seven-point construction using symbolic arithmetic.
A separate numerical stress test checks the common diagonal surrogate
for 1000 Hermitian examples of dimensions two through six. This
numerical check supplements the analytic proof and is not a substitute
for it.
Numerical searches helped locate the witness; the theorems above
do not rely on a finite grid, floating-point positivity, a fitted
polynomial, or a semidefinite solver's status.

The concrete advance proposed here is a finite-rate and quantitatively
sharp compatibility test. Earlier Hamiltonian-ensemble work supplies
the conceptual question, and classical rational quadrature explains
the contrasting qubit aliases. The paper does not solve general
quantum process identification, establish a microscopic source of
noise, or assert a universal boundary for all dephasing generators.
There is no experimental validation in this manuscript. Noise with a
finite correlation time is a separate interpolation problem: rejection
of a static law alone does not estimate that correlation time or identify
an Ornstein--Uhlenbeck process. Echo-based methods and non-random-unitary
dephasing tests \cite{hahn,helm} address different control or model classes.

Natural next problems are to optimize over other rate ratios and
measurement budgets, to determine the sharp distance for the full
two-state data, and to sharpen the population-calibration constants. Higher-dimensional energy-difference constraints may supply
further finite resolvent witnesses. These are open directions, not
results claimed here.

\paragraph{AI assistance and reproducibility.}
AI systems assisted the literature search, exploratory computation,
proof development and writing. A GPT-6 Astra advisory review contributed
literature suggestions and the spectral-monotonicity sharpening of the
population bound; it is not a blind independent review. A separate
Claude Opus 5.5 review of the preceding draft found no fatal error but
rated its significance substantially lower. Both reports and a response
are preserved with the source; the present editorial revision has not
received a new score. Reviews are advisory and are not
formal proof certificates or evidence of scientific priority.
The companion source and verification script make every algebraic
claim in the main witness auditable. The manuscript is an unsubmitted
candidate; no publication decision or external peer-review outcome
is implied.

\begin{thebibliography}{99}
\bibitem{kropf} C. M. Kropf, C. Gneiting and A. Buchleitner,
Effective dynamics of disordered quantum systems,
\emph{Phys. Rev. X} \textbf{6}, 031023 (2016).
\href{https://doi.org/10.1103/PhysRevX.6.031023}{doi:10.1103/PhysRevX.6.031023}.
\bibitem{chen2018} H.-B. Chen, C. Gneiting, P.-Y. Lo, Y.-N. Chen
and F. Nori, Simulating open quantum systems with Hamiltonian
ensembles and the nonclassicality of the dynamics,
\emph{Phys. Rev. Lett.} \textbf{120}, 030403 (2018).
\href{https://doi.org/10.1103/PhysRevLett.120.030403}{doi:10.1103/PhysRevLett.120.030403}.
\bibitem{chen2019} H.-B. Chen, P.-Y. Lo, C. Gneiting, J. Bae,
Y.-N. Chen and F. Nori, Quantifying the nonclassicality of pure
dephasing, \emph{Nature Communications} \textbf{10}, 3794 (2019).
\href{https://doi.org/10.1038/s41467-019-11502-4}{doi:10.1038/s41467-019-11502-4}.
\bibitem{chen2021} H.-B. Chen and Y.-N. Chen, Canonical Hamiltonian
ensemble representation of dephasing dynamics and the impact of
thermal fluctuations on quantum-to-classical transition,
\emph{Scientific Reports} \textbf{11}, 10046 (2021).
\href{https://doi.org/10.1038/s41598-021-89400-3}{doi:10.1038/s41598-021-89400-3}.
\bibitem{rose} D. C. Rose, H. Touchette, I. Lesanovsky and J. P.
Garrahan, Spectral properties of simple classical and quantum reset
processes, \emph{Phys. Rev. E} \textbf{98}, 022129 (2018).
\href{https://doi.org/10.1103/PhysRevE.98.022129}{doi:10.1103/PhysRevE.98.022129}.
\bibitem{lubinsky} D. S. Lubinsky, Explicit orthogonal polynomials for
reciprocal polynomial weights, \emph{Proc. Amer. Math. Soc.}
\textbf{137}, 2317--2327 (2009).
\href{https://lubinsky.math.gatech.edu/Research%20papers/ExplicitOPAugust2008.pdf}{Author manuscript}.
\bibitem{wang} B. Wang and R. Wang, Joint mixability,
\emph{Mathematics of Operations Research} \textbf{41}, 808--826 (2016).
\href{https://doi.org/10.1287/moor.2015.0755}{doi:10.1287/moor.2015.0755}.
\bibitem{fagnola} F. Fagnola, J. E. Gough, H. I. Nurdin and L. Viola,
Mathematical models of Markovian dephasing (2019).
\href{https://arxiv.org/abs/1811.11784}{arXiv:1811.11784}.
\bibitem{schirmer} S. G. Schirmer and A. I. Solomon,
Constraints on relaxation rates for $N$-level quantum systems,
\emph{Phys. Rev. A} \textbf{70}, 022107 (2004).
\href{https://doi.org/10.1103/PhysRevA.70.022107}{doi:10.1103/PhysRevA.70.022107}.
\bibitem{lin} J.-D. Lin et al., Resource efficient certification of
system environment entanglement solely from reduced system dynamics
(2025). \href{https://arxiv.org/abs/2510.17140}{arXiv:2510.17140}.
\bibitem{eriksson-maintenance} L. Eriksson, Maintenance is not
restoration: endpoint no-go theorems, exact SSH resource horizons,
and quantum rapid-replacement equality (2026).
\href{https://www.ai.vixra.org/abs/2607.0031}{ai.viXra:2607.0031}.
\bibitem{eriksson-moments} L. Eriksson, Finite-sample-valid tests for
structured matricial Hausdorff moments: joint Gaussian Grams, strict
half-time separation, and sharp tangent-cone power (2026).
\href{https://airr.science/papers/ARR-2026-7X5XX0CBE19MXVSZ/}{AIRR ARR-2026-7X5XX0CBE19MXVSZ}.

\bibitem{hahn} E. L. Hahn, Spin echoes,
\emph{Phys. Rev.} \textbf{80}, 580--594 (1950).
\href{https://doi.org/10.1103/PhysRev.80.580}{doi:10.1103/PhysRev.80.580}.
\bibitem{kummerer} B. K\"ummerer and H. Maassen,
The essentially commutative dilations of dynamical semigroups on $M_n$,
\emph{Commun. Math. Phys.} \textbf{109}, 1--22 (1987).
\href{https://doi.org/10.1007/BF01205670}{doi:10.1007/BF01205670}.
\bibitem{kemperman} J. H. B. Kemperman,
The general moment problem, a geometric approach,
\emph{Ann. Math. Statist.} \textbf{39}, 93--122 (1968).
\href{https://doi.org/10.1214/aoms/1177698508}{doi:10.1214/aoms/1177698508}.
\bibitem{bayer} C. Bayer and J. Teichmann,
The proof of Tchakaloff's theorem,
\emph{Proc. Amer. Math. Soc.} \textbf{134}, 3035--3040 (2006).
\href{https://arxiv.org/abs/math/0502473}{arXiv:math/0502473}.
\bibitem{agler} J. Agler and J. E. McCarthy,
The use of kernel functions in solving the Pick interpolation problem
(2014). \href{https://www.math.wustl.edu/~mccarthy/public_papers/KernelsandPick.pdf}{Author manuscript}.
\bibitem{puccetti} G. Puccetti, P. Rigo, B. Wang and R. Wang,
Centers of probability measures without the mean (2018 manuscript).
\href{https://arxiv.org/abs/1704.02660}{arXiv:1704.02660}.
\bibitem{helm} J. Helm and W. T. Strunz,
Quantum decoherence of two qubits,
\emph{Phys. Rev. A} \textbf{80}, 042108 (2009).
\href{https://arxiv.org/abs/0910.5364}{arXiv:0910.5364}.
\end{thebibliography}
\end{document}
